\documentclass[11pt,a4paper,UTF8,fontset=none]{ctexart}
\setCJKmainfont{Noto Serif CJK SC}
\setCJKsansfont{Noto Sans CJK SC}
\setmainfont{DejaVu Serif}
\usepackage[a4paper,margin=22mm]{geometry}
\usepackage{amsmath,amssymb,booktabs,longtable,array,tabularx,xcolor,hyperref,fancyhdr,enumitem,microtype}
\hypersetup{colorlinks=true,linkcolor=blue,urlcolor=blue}
\definecolor{navy}{HTML}{173452}
\definecolor{sky}{HTML}{E9F4FA}
\pagestyle{fancy}\fancyhf{}\fancyhead[L]{\sffamily ZTQI · 研究试验版 v0.2}\fancyhead[R]{\sffamily 2026-10-21}\fancyfoot[C]{\thepage}
\setlength{\parindent}{2em}\setlength{\parskip}{4pt}
\begin{document}
\begin{center}{\Huge\sffamily\bfseries ZTQI-72 研究问卷}\par\vspace{2mm}{\large 纸质版 v0.2 · 2026-10-21}\end{center}
\noindent\fcolorbox{black}{white}{%
  \parbox{\dimexpr\linewidth-2\fboxsep-2\fboxrule\relax}{%
    \textbf{自愿填写说明}：本问卷是未经验证的研究原型，不用于判断他人的道德、人格或心理健康。全部 72 道题均请作答，预计 12--18 分钟。无需填写姓名。请按最近三个月的真实想法选择。对于标注\mbox{“过去一个月”}的条目请按其指定时间作答。答题卡须用 2B 铅笔填满对应椭圆，切勿折叠。%
  }%
}\par\vspace{2mm}
\noindent\textbf{统一选项（全部 72 道题通用）}\par\smallskip
\noindent\begin{tabularx}{\linewidth}{@{}*{5}{>{\centering\arraybackslash}X}@{}}
\textbf{1} 非常不同意 & \textbf{2} 比较不同意 & \textbf{3} 中立/不确定 & \textbf{4} 比较同意 & \textbf{5} 非常同意
\end{tabularx}\par\vspace{3mm}
\begin{longtable}{@{}p{1.2cm}p{14.5cm}@{}}
\toprule 题号&陈述（在独立答题卡上填涂）\\\midrule\endhead
01&我投入很多时间学习后，会觉得别人理应给予我额外的尊重\\[3pt]
02&我从来没有因为嫉妒而说过任何违心的话\\[3pt]
03&成绩较差的人也可能在重要问题上比我判断得更准确\\[3pt]
04&我愿意接受一些重要的人生问题没有最优解\\[3pt]
05&别人意见与我不同时，我常把分歧看作有人没有理解标准答案\\[3pt]
06&我觉得学业成绩可以作为评价一个人整体能力的主要标准\\[3pt]
07&过去一个月，我参与过与考试成绩没有直接关系的兴趣活动\\[3pt]
08&我付出的学习努力应当换来同学更多的关注\\[3pt]
09&我担心成绩平庸会使我变成一个没有价值的人\\[3pt]
10&过去一个月，我曾主动了解成绩不如我的同学擅长的领域\\[3pt]
11&我有时觉得学业优秀应当抵消我在人际交往中的不足\\[3pt]
12&被拒绝时，我能够尊重对方的自主选择而不以成绩辩解\\[3pt]
13&我不认为成绩优势能赋予我更多的交友或恋爱资格\\[3pt]
14&当自己的优秀没有得到认可时，我容易认为是别人不识货\\[3pt]
15&努力学习并不能使别人负有欣赏我的义务\\[3pt]
16&成绩不好的同学提出观点时，我会先怀疑其判断能力\\[3pt]
17&对于价值观冲突，我常试图证明一方必然比另一方正确\\[3pt]
18&当成绩不如我的人更受欢迎时，我会觉得秩序颠倒了\\[3pt]
19&只知道两个人成绩时，我会倾向于认为成绩更好的人整体上更优秀\\[3pt]
20&一个人的成绩与他应得的基本尊重没有关系\\[3pt]
21&理解不同人的经历，有时比给出答案更重要\\[3pt]
22&我不会根据一次考试成绩推断一个人的品格\\[3pt]
23&别人有权不喜欢我，即使我的成绩比大多数人更好\\[3pt]
24&我的自信很大程度上来自考试排名\\[3pt]
25&不同领域的努力都可能产生值得欣赏的结果\\[3pt]
26&我从未有过任何不公平的想法\\[3pt]
27&只要掌握足够多的规则，我觉得人际交往就能像解题一样稳定\\[3pt]
28&如果一个人的成绩差，我会觉得他在其他领域的成功价值有限\\[3pt]
29&如果未来不再擅长考试，我会觉得自己的身份失去了支撑\\[3pt]
30&一个不擅长考试的人，也完全可能获得有意义的成就\\[3pt]
31&遇到人际冲突时，我首先想找出唯一正确的一方\\[3pt]
32&成绩比我差很多的人对我的建议，我通常不会太在意\\[3pt]
33&我能真诚欣赏成绩远不如我的人的其他长处\\[3pt]
34&为核对问卷填写质量，请在本题选择“比较不同意”（2）\\[3pt]
35&当一个人成绩好时，我会更容易原谅他在人际相处中的缺点\\[3pt]
36&即使其他方面进步很大，考试失利仍会让我觉得自己一无是处\\[3pt]
37&我常把相处中的挫折归因于自己没找到正确解法\\[3pt]
38&若无法提升成绩，我认为多数课外活动不值得投入\\[3pt]
39&相比于考试高分，我通常不太看重一个人的创造力\\[3pt]
40&我认为将大量精力放在非应试活动上，大多是在浪费潜力\\[3pt]
41&我永远不会为了表现好而隐瞒自己的真实想法\\[3pt]
42&即使无法明确判定对错，讨论仍然可能富有价值\\[3pt]
43&我是否值得被尊重，并不取决于最近一次考试\\[3pt]
44&我可以接受自己在某些学科上并不出色\\[3pt]
45&在评价一个陌生人时，我通常把学习成绩放在最重要的位置\\[3pt]
46&我很难理解有人自愿选择不以学业竞争为中心的生活\\[3pt]
47&在学校里，成绩越优秀的人一般越有资格获得话语权\\[3pt]
48&没有优异成绩，我也能坦然地向他人介绍自己\\[3pt]
49&我经常用分数来判断自己今天过得是否有意义\\[3pt]
50&我认为学业表现好的人在人际关系中理应有某些优先权\\[3pt]
51&为核对问卷填写质量，请在本题选择“比较同意”（4）\\[3pt]
52&别人称赞我成绩好时，我才比较确信自己足够优秀\\[3pt]
53&尊重他人的生活选择并不要求其符合我的升学目标\\[3pt]
54&过去一个月，我曾因考试成绩影响对某位同学的态度\\[3pt]
55&我认为大多数复杂的人际问题都可以归结为一套正确规则\\[3pt]
56&人际分歧有时确实可以同时存在多个合理立场\\[3pt]
57&复杂情境中，我通常不喜欢接受含糊不清的结论\\[3pt]
58&即使考试失利，我仍能清楚说出自己值得肯定的特点\\[3pt]
59&过去一个月，我曾与他人讨论没有唯一答案的价值问题\\[3pt]
60&考试排名往往可以反映一个人是否值得敬佩\\[3pt]
61&在我看来，学习成绩一般比诚实、善良等品质更能说明人的价值\\[3pt]
62&当成绩下降时，我很难保持对自己的基本认可\\[3pt]
63&如果我成绩突出却没有得到某个人的好感，我会觉得不太公平\\[3pt]
64&过去一个月，我曾因考试失利回避与人交往\\[3pt]
65&我在所有情境下都会完全客观地评价别人\\[3pt]
66&同等条件下，我认为成绩更好的人更应得到朋友的偏爱\\[3pt]
67&面对社会问题，我经常只关注是否能找到明确的对错判断\\[3pt]
68&没有升学优势的兴趣爱好也可以是人生的重要组成部分\\[3pt]
69&我认为艺术或体育方面的成就通常不如高考试成绩值得尊重\\[3pt]
70&一次重要考试不理想会让我怀疑自己作为一个人的价值\\[3pt]
71&花很多时间发展与升学无关的兴趣，通常是不明智的\\[3pt]
72&过去一个月，我曾因为他人的拒绝而联想到自己的成绩\\[3pt]
\bottomrule\end{longtable}
\vspace{3mm}\noindent\textbf{结束确认}：请检查答题卡 01--72 是否一一作答。如不愿提交，可直接收回或销毁答题卡。\end{document}